\introduction % Структурный элемент: ВВЕДЕНИЕ

В последние годы технологии автономного вождения перешли из стадии лабораторных прототипов к практическому применению в городских транспортных сервисах. Наиболее заметным примером являются беспилотные такси Waymo~\cite{waymo-habr}, эксплуатируемые на дорогах общего пользования и демонстрирующие возможность длительной автономной работы в реальной городской среде. В России развитие беспилотного транспорта также связано с проектом <<Автономный транспорт Яндекса>>~\cite{yandex-autonomy}, в рамках которого разрабатываются роботакси~\cite{yandex-robotaxi}, автономные грузовики~\cite{yandex-truck} и роботы-доставщики~\cite{yandex-rovers}. Подобные системы функционируют в динамической и частично наблюдаемой среде, в которой качество оценки собственного движения напрямую влияет на безопасность планирования и управления.

Одним из ключевых источников информации о геометрии окружающего пространства является лидар, обеспечивающий измерения расстояний и позволяющий восстанавливать структуру сцены даже при слабой текстуре и изменении освещённости. Поэтому лидарная одометрия остаётся важным компонентом систем локализации беспилотного транспорта, особенно в сценариях, где использование только визуальных методов оказывается недостаточно надёжным. Дополнительную значимость таким методам придаёт и то, что спутниковая навигация по сигналам GNSS в реальных условиях может работать нестабильно. В частности, в России качество такого сигнала может значительно ухудшаться, что повышает требования к автономным методам оценки собственного движения. Вместе с тем практическое применение лидарной одометрии связано с рядом ограничений.

Существенную проблему для алгоритмов локализации представляют крупные динамические объекты, например грузовые автомобили, которые могут перекрывать значительную часть поля зрения лидара. В такой ситуации уменьшается доля устойчивых геометрических ориентиров, ухудшается качество поиска соответствий между облаками точек и возрастает риск смещения оценки относительного перемещения. Для беспилотного автомобиля это означает деградацию локализации и снижение качества траекторного планирования.

В рассматриваемой работе интерес представляет практический сценарий, в котором при временном перекрытии обзора грузовиком необходимо сохранить достаточно точную оценку движения, чтобы беспилотный автомобиль мог надёжно проехать небольшой участок пути, опираясь преимущественно на одометрическую оценку. Речь не идёт о длительном автономном движении только по одометрии; задача состоит в том, чтобы разработать алгоритм лидарной одометрии, учитывающий перекрытие поля зрения и обеспечивающий устойчивую работу на коротком интервале времени до восстановления большей наблюдаемости сцены. Именно такие кратковременные, но критически важные эпизоды представляют особый интерес для прикладных систем автономного вождения.

Актуальность работы определяется необходимостью повышения устойчивости алгоритмов лидарной одометрии в условиях частичного перекрытия поля зрения динамическими объектами. Несмотря на большое количество исследований в области регистрации облаков точек, задача сохранения точности одометрии при длительном перекрытии значимой части сцены остаётся практически важной и недостаточно тривиальной. Особенно это касается сценариев движения в транспортном потоке, когда грузовик, движущийся рядом или выполняющий обгон, временно закрывает существенную часть наблюдаемого пространства и тем самым искажает структуру входных данных.

Объектом исследования являются методы локализации и оценки движения беспилотного автомобиля на основе лидарных данных. Предметом исследования являются алгоритмы лидарной одометрии, учитывающие перекрытие поля зрения при регистрации облаков точек и фильтрации оценок движения.

Целью диссертационной работы является разработка устойчивого алгоритма лидарной одометрии беспилотного автомобиля, учитывающего перекрытие поля зрения и позволяющего оценивать движение на коротком участке пути.

Для достижения поставленной цели в работе решаются следующие задачи:
\begin{enumerate}
    \item провести анализ существующих подходов к лидарной одометрии и выявить ограничения классических методов регистрации облаков точек в условиях динамического перекрытия поля зрения;
    \item описать используемый набор данных Boreas-RT~\cite{boreas-road-trip} и определить, какие данные из него применяются в экспериментальном исследовании;
    \item формализовать сценарий перекрытия поля зрения лидара и определить требования к устойчивости алгоритма оценки движения на коротком участке пути;
    \item разработать алгоритм обработки последовательности лидарных кадров, включающий вокселизацию облаков точек, регистрацию двух соседних облаков точек, оценку преобразования между парами точек с учётом весов и применение модели движения;
    \item определить метрики качества и реализовать программный конвейер для запуска серии вычислительных экспериментов на наборах лидарных данных с возможностью моделирования перекрытия поля зрения;
    \item провести экспериментальное исследование точности алгоритма по компонентам линейной и угловой скорости, а также определить, на каких данных и при каких параметрах предложенный подход работает наиболее устойчиво.
\end{enumerate}

В работе исследуется устойчивый алгоритм, в котором оценка относительного движения между последовательными лидарными кадрами строится на основе итеративной регистрации двух соседних облаков точек. Для повышения устойчивости используется оценка преобразования между парами точек с учётом весов, а для учёта кинематической инерции движения --- начальное приближение, формируемое по оценённым скоростям, и фильтр Калмана~\cite{habr-kalman-pictures} с моделью постоянной скорости. Такой подход обеспечивает снижение влияния выбросов, повышает устойчивость оценки движения в условиях неполной наблюдаемости сцены и позволяет пройти короткий участок траектории без существенной деградации локализации.

Методы исследования включают методы трёхмерной геометрической регистрации, численные методы оптимизации, методы статистической обработки данных и вычислительный эксперимент. В качестве базового метода оценки относительного преобразования используется алгоритм ICP~\cite{icp-learnopencv} с поиском ближайших соответствий между двумя соседними облаками точек и вычислением евклидова преобразования по алгоритму Кабша~\cite{kabsch-wikipedia,kabsch-nghiaho}. Для подавления влияния выбросов оценка преобразования между парами точек выполняется с учётом весов, дающих устойчивость к выбросам, в частности с использованием функции Гемана--Макклюра~\cite{kiss-icp}. Сглаживание и предсказание скоростей выполняются с помощью фильтра Калмана~\cite{habr-kalman-pictures} в предположении модели постоянной скорости. Экспериментальное исследование проводится на данных набора Boreas-RT~\cite{boreas-road-trip} с использованием эталонных значений линейных и угловых скоростей; качество работы алгоритма оценивается по метрикам среднего смещения, стандартного отклонения, MSE, MAE, RMSE и 95\%-квантиля абсолютной ошибки. Для анализа распределения ошибок используются гистограммы, для визуального сопоставления оценённых и эталонных скоростей --- временные графики, а для сравнения различных конфигураций алгоритма --- парный критерий Стьюдента~\cite{student-yandex-practicum,student-t-test-habr}, применяемый к метрикам, рассчитанным на одинаковых отрезках проезда. Таким образом, исследование сочетает алгоритмический анализ, программную реализацию и статистически обоснованную проверку полученных результатов.

Научная новизна работы состоит в разработке и исследовании устойчивого алгоритма лидарной одометрии, учитывающего перекрытие поля зрения, а также в построении экспериментального конвейера, позволяющего сравнивать различные конфигурации алгоритма в воспроизводимых условиях. Практическая значимость работы заключается в возможности использования предложенного подхода для повышения надёжности оценки движения беспилотного автомобиля в дорожных сценах, когда требуется устойчиво пройти короткий участок пути до восстановления большей наблюдаемости.

Структурно диссертация состоит из введения, трёх разделов основной части, заключения, списка использованных источников и приложений. В первом разделе рассматриваются особенности технологии лидарной одометрии, алгоритм ICP и влияние динамического перекрытия поля зрения на качество оценки движения. Во втором разделе описываются набор данных Boreas-RT~\cite{boreas-road-trip}, этапы предобработки облаков точек, алгоритм ICP с оценкой преобразования между парами точек с учётом весов, модель движения, фильтр Калмана и экспериментальный конвейер с используемыми метриками качества. В третьем разделе приводятся результаты вычислительных экспериментов, анализируется устойчивость алгоритма и обсуждаются полученные зависимости. В заключении формулируются основные результаты работы, а также отмечаются ситуации, в которых предложенный подход показал устойчивость, и случаи, в которых его возможности остаются ограниченными.
